% ============================================================================
% MODULE 2: FIRST SIMULATION
% ============================================================================
% Duration: 60 minutes (hands-on; near-solved file provided)
% Learning Goals: Set up, run, and interpret a fallow water balance simulation
% Prerequisites: Module 0, Module 1

\chapter{Module 2: First Simulation}
\label{ch:module2}

\section{Learning Objectives}

By the end of this module, you will:
\begin{itemize}
  \item Set up a simple APSIM simulation from scratch
  \item Understand fallow water balance: rainfall, runoff, drainage, evaporation
  \item Interpret daily ESW (extractable soil water) as an indicator of sowing readiness
  \item Run simulation and export results to Excel
  \item Troubleshoot common startup errors (file not found, missing dates, etc.)
\end{itemize}

\section{Conceptual Framework: Fallow Water Balance}

\subsection{What is a Fallow?}

A fallow is a season where the field is left without a crop. Soil water accumulates (useful before sowing) or is lost to drainage/evaporation.

\begin{keyconceptbox}
\textbf{Daily Water Balance Equation:}
\[
\text{ESW}_{t} = \text{ESW}_{t-1} + \text{Rainfall} - \text{Runoff} - \text{Drainage} - \text{Evaporation}
\]

where ESW = extractable soil water (mm), available for plant growth.
\end{keyconceptbox}

\subsection{When Are Soils Ready for Sowing?}

Farmers typically wait for:
\begin{itemize}
  \item ESW $>$ 100 mm (soil has adequate water to support germination and early growth)
  \item Rainfall $>$ 25 mm in last 7 days (recent rain indicates favorable conditions)
  \item Date within sowing window (e.g., May 1 – July 10 for winter wheat in Queensland)
\end{itemize}

Exercise A explores these thresholds using a 7-year fallow simulation.

\section{Getting Started: Open Exercise A}

\subsection{File Location}

\filename{Exercise\_A\_Fallow\_Water\_Balance.apsimx} (provided in course materials)

\subsection{Simulation Setup}

\textbf{Clock:} 7 years of a representative site (e.g., 1950–1957)

\textbf{Weather:} Dalby, Queensland (real BOM/SILO data)

\textbf{Soil:} Norwin Vertosol (black clay; realistic for wheat region)

\textbf{Management:} None (fallow = no crop, no fertiliser)

\textbf{Report Variables:} 
\begin{itemize}
  \item [Clock].Today (date)
  \item [Soil].ESW (extractable soil water)
  \item [Weather].Rainfall (daily rain)
  \item [Soil].SoilWater.Evaporation (simulated evaporation)
  \item [Soil].SoilWater.Drainage (deep percolation)
\end{itemize}

\section{Running the Simulation}

\subsection{Step 1: Open APSIM}

\begin{enumerate}
  \item Launch APSIM Next Gen
  \item File $\to$ Open $\to$ select \filename{Exercise\_A\_Fallow\_Water\_Balance.apsimx}
\end{enumerate}

\subsection{Step 2: View Structure}

\begin{itemize}
  \item Expand the tree on the left
  \item Verify components: Clock, Weather, Soil, Report, DataStore
  \item Click each component to check configuration (dates, file paths, variables)
\end{itemize}

\subsection{Step 3: Run Simulation}

\begin{enumerate}
  \item Click the Simulation in the tree
  \item Press Run button (or Ctrl+R)
  \item Watch the output window for progress
  \item Simulation should complete without errors
\end{enumerate}

\subsection{Step 4: View Results}

\begin{enumerate}
  \item Click DataStore in the tree
  \item Right-click $\to$ Export $\to$ To Excel
  \item Name file: \filename{Exercise\_A\_Results.xlsx}
  \item Open in Excel; examine data
\end{enumerate}

\section{Analyzing the Data}

\subsection{Key Columns}

\begin{table}[H]
\centering
\caption{Output Variables in Exercise A}
\label{tab:exercise-a-vars}
\begin{tabularx}{\textwidth}{l|X|X}
\toprule
\textbf{Column} & \textbf{Units} & \textbf{Interpretation} \\
\midrule
Clock.Today & Date & Calendar date \\
Soil.ESW & mm & Plant-available water; sowing threshold is ESW $>$ 100 mm \\
Weather.Rainfall & mm & Daily precipitation \\
Soil.SoilWater.Evaporation & mm & Soil water lost to atmosphere (PET if no crop) \\
Soil.SoilWater.Drainage & mm & Water percolating below root zone \\
\bottomrule
\end{tabularx}
\end{table}

\subsection{Analysis Prompts}

\begin{enumerate}
  \item \textbf{Seasonal Pattern:} Plot ESW vs. date across all 7 years. When does ESW peak? When is it lowest?
  
  \item \textbf{Sowing Window:} Identify dates when ESW $>$ 100 mm AND rainfall in last 7 days $>$ 25 mm. How many sowing opportunities per year?
  
  \item \textbf{Dry Year vs. Wet Year:} Compare ESW profiles for a dry year vs. wet year. What drives the difference?
  
  \item \textbf{Evaporation vs. Drainage:} Which process removes more water? Does it vary seasonally?
\end{enumerate}

\section{Common Errors and Troubleshooting}

\begin{warningbox}
\textbf{Error 1: ``File not found'' for weather file}

\textit{Solution:} Weather file path uses \texttt{\%root\%} variable (APSIM's installation directory). Verify weather file exists at: \texttt{\%APSIM\_Installation\%/Examples/WeatherFiles/AU\_Dalby.met}. If not found, check APSIM installation or download file separately.
\end{warningbox}

\begin{warningbox}
\textbf{Error 2: ``Soil profile empty'' or ``missing layers''

\textit{Solution:} Ensure Soil component has Physical, SoilWater, and Nutrient sub-components. These are auto-created if you load a predefined soil from APSOIL database. If manually editing \filename{.apsimx}, verify JSON structure is correct.
\end{warningbox}

\begin{warningbox}
\textbf{Error 3: ``Report variables not found''

\textit{Solution:} Verify variable names match APSIM's internal paths. Example: \texttt{[Soil].ESW} not \texttt{Soil.ESW}; \texttt{[Weather].Rainfall} not \texttt{Weather.Rain}. Consult APSIM GUI (right-click Report) to auto-generate valid variable names.
\end{warningbox}

\section{Extending the Analysis}

\subsection{Comparison: Different Soils}

Create two variants of Exercise A:
\begin{enumerate}
  \item Same as above (black clay)
  \item Shallow sandy soil (low water-holding capacity)
\end{enumerate}

Run both; compare ESW profiles. Which soil reaches sowing threshold more often?

\subsection{Comparison: Different Weather}

Use a different weather file (e.g., wetter/drier region). How does ESW dynamics change?

\subsection{Multi-Year Statistics}

Calculate for all 7 years:
\begin{itemize}
  \item Average ESW by month
  \item Frequency of sowing-ready days
  \item Coefficient of variation (year-to-year variability)
\end{itemize}

\section{Conceptual Questions}

\begin{enumerate}
  \item Why is ESW in the fallow higher in spring than summer, even if rainfall is similar?
  
  \item If you sow wheat on day X when ESW = 110 mm, does ESW stay above 100 mm? Why or why not?
  
  \item How would deep drainage change if the soil were twice as thick (double the LL and DUL)?
  
  \item Could a rainy day during grain fill harm wheat yield? How would APSIM capture this?
\end{enumerate}

\section{Summary}

\begin{itemize}
  \item Fallow water balance is foundational: daily rain/runoff/drainage/evaporation
  \item Sowing readiness depends on ESW and recent rain (decision threshold)
  \item Exercise A provides 7-year dataset for analyzing water availability
  \item Results can be exported to Excel for custom analysis and visualization
  \item Comparing soil types and weather regions reveals biophysical constraints
\end{itemize}

\section{Key Resources}

\begin{itemize}
  \item [[ESW\_Extractable\_Soil\_Water]] — Definition and role in sowing
  \item [[Water\_Balance]] — Daily water accounting
  \item [[Soil]] — Soil profile configuration
  \item [[Manager]] — Sowing logic (for future reference)
\end{itemize}

\section{What's Next?}

In Module 3, you'll learn to interpret simulation outputs: calculating averages, plotting trends, and extracting agronomic insights from data tables.

% ============================================================================
% END MODULE 2
% ============================================================================
